\documentclass[12pt, a4paper, oneside]{report}
%==================================================
% IMPORTS
% Package utama untuk layout, tabel, gambar, dan utilitas lainnya
%==================================================

\usepackage{tocloft}            % Custom daftar isi, gambar, tabel, referensi
\usepackage{tabularx}           % Tabel dengan lebar kolom fleksibel
\usepackage{longtable}          % Tabel multi-halaman
\usepackage{multirow}           % Merge baris dalam tabel
\usepackage{multicol}           % Layout multi-kolom
\usepackage{array}              % Kontrol format tabel lebih lanjut
\usepackage{booktabs}           % Garis tabel lebih clean (toprule, midrule, bottomrule)

\usepackage{graphicx}           % Insert gambar
\usepackage[export]{adjustbox}  % Styling tambahan untuk gambar (misalnya border)

\usepackage{ragged2e}           % Alignment teks (lebih fleksibel dari default)
\usepackage{tabto}              % Kontrol indentasi manual

\usepackage{enumitem}           % Custom bullet & numbering list

\usepackage{listings}           % Menampilkan code
\usepackage{xcolor}             % Warna (dipakai oleh listings)

\usepackage{pdfpages}           % Include file PDF
\usepackage{pgffor}             % Looping di LaTeX

\usepackage{float}              % Kontrol posisi figure/table (H, h, dll)
\usepackage{lipsum}             % Dummy text (opsional)

\usepackage[indonesian]{babel}  % Bahasa dokumen

%==================================================
% METADATA
%==================================================

\title{
    {\fontsize{22pt}{40pt}\selectfont
            MAGANG INDUSTRI SEBAGAI \textit{FRONT-END DEVELOPER} DI PT INTI UTAMA SOLUSINDO}
}

\newcommand{\documenttitle}{
    \MakeUppercase{
        \jenisMBKM\@
        SEBAGAI\\
        \textit{\jobposition}\\
        DI\@ \company\\
    }
}

\newcommand{\smalldocumenttitle}{\jenisMBKM\ sebagai \jobposition\ di \company}

% Biodata
\newcommand{\authorname}{Aldian Yohanes}
\newcommand{\authornim}{535230139}
\newcommand{\studyprogram}{Teknik Informatika}
\newcommand{\faculty}{Teknologi Informasi}
\newcommand{\semester}{GENAP}
\newcommand{\academicperiod}{2025-2026}

\newcommand{\jobposition}{Full-Stack Developer}
\newcommand{\jobdepartment}{NPD Pharmalogic / Hospinet Platform}
\newcommand{\company}{PT Inti Utama Solusindo}
% Data MBKM
\newcommand{\jenisMBKM}{Magang Industri}
\newcommand{\lamaMBKM}{12 bulan}
\newcommand{\mentor}{Dimas Aditya Nugroho}
\newcommand{\dosen}{Viny Christanti Mawardi, M.Kom.}
\newcommand{\koordinator}{Tony, Ph.D.}
\newcommand{\city}{Jakarta}
\newcommand{\signaturedate}{9 Juni 2026}
\newcommand{\reportmonth}{Juni}
\newcommand{\reportyear}{2026}
% Data Dokumen
\newcommand{\jenisdokumen}{LAPORAN AKHIR }   % Data seperti judul, penulis, dll

%==================================================
% FONT CONFIGURATION (XeLaTeX)
%==================================================

\usepackage{fontspec}
\setmainfont{Calibri}           % Font utama dokumen
\setmonofont{Consolas}[LetterSpace=-5.0] % Font untuk code

\usepackage{microtype}          % Optimasi spacing teks (lebih rapi saat justify)

%==================================================
% PAGE LAYOUT (MARGIN)
%==================================================

\usepackage[
    a4paper,
    left=4cm,
    right=3cm,
    top=4cm,
    bottom=3cm,
    % showframe % Aktifkan jika ingin debug margin
]{geometry}

%==================================================
% HEADER & FOOTER
%==================================================

\usepackage{fancyhdr}

\pagestyle{fancy}
\fancyhf{}                      % Reset semua header & footer

\renewcommand{\headrulewidth}{0pt} % Hilangkan garis header
\renewcommand{\footrulewidth}{0pt} % Hilangkan garis footer

\fancyfoot[R]{\thepage}         % Nomor halaman di kanan bawah

% Style khusus untuk halaman "plain" (misalnya awal bab)
\fancypagestyle{plain}{
    \fancyhf{}
    \renewcommand{\headrulewidth}{0pt}
    \renewcommand{\footrulewidth}{0pt}
    \fancyfoot[R]{\thepage}
}

%==================================================
% SPACING & INDENTASI
%==================================================

\usepackage{setspace}

\onehalfspacing%               % Line spacing 1.5
\setlength{\parindent}{1.27cm} % Indent paragraf pertama

\usepackage{indentfirst}        % Paragraf pertama tetap indent

%==================================================
% HEADING FORMAT
%==================================================

\usepackage{titlesec}

% ---------- BAB (Heading 1) ----------
\titleformat{\chapter}[display]
{\normalfont\bfseries\fontsize{14pt}{16pt}\centering}
{BAB \thechapter}
{0pt}
{\vspace{0.5ex}}

\titlespacing*{\chapter}
{0pt}
{-20pt}
{20pt}

% ---------- Subbab (Heading 2) ----------
\titleformat{\section}
{\normalfont\bfseries\fontsize{12pt}{14pt}\selectfont}
{\thesection}
{10pt}
{}

% ---------- Subsubbab (Heading 3) ----------
\titleformat{\subsection}
{\normalfont\fontsize{12pt}{14pt}\selectfont}
{\thesubsection}
{10pt}
{}

\titlespacing{\subsection}{0pt}{12pt}{0pt}

%==================================================
% TEXT FORMATTING (ANTI HYPHEN)
%==================================================

% Mengurangi pemotongan kata di akhir baris
\hyphenpenalty=5000
\exhyphenpenalty=10000
\tolerance=9999
\emergencystretch=10pt
\hbadness=99999

%==================================================
% HYPERLINK
%==================================================

\usepackage{hyperref}

\hypersetup{
    colorlinks=true,
    linkcolor=black,
    filecolor=black,
    urlcolor=blue,
    citecolor=black
}

%==================================================
% CAPTION (GAMBAR & TABEL)
%==================================================

\usepackage{caption}

\captionsetup[figure]{
    labelsep=space,
    justification=centering
}

%==================================================
% LABEL BAHASA INDONESIA
%==================================================

\addto\captionsindonesian{
    \renewcommand{\contentsname}{DAFTAR ISI}
    \renewcommand{\refname}{Referensi}
    \renewcommand{\listfigurename}{DAFTAR GAMBAR}
    \renewcommand{\listtablename}{DAFTAR TABEL}
    \renewcommand{\lstlistlistingname}{DAFTAR KODE}
    \renewcommand{\figurename}{Gambar}
    \renewcommand{\tablename}{Tabel}
    \renewcommand{\lstlistingname}{Kode}
    \renewcommand{\bibname}{DAFTAR PUSTAKA}
}

%==================================================
% CODE LISTING STYLE
%==================================================

\lstset{
    basicstyle=\ttfamily\small,
    frame=single,
    breaklines=true,
    tabsize=2,
    numbers=left,
    numberstyle=\tiny,
    keywordstyle=\color{blue!70!black},
    commentstyle=\color{gray},
    stringstyle=\color{red!70!black},
    showstringspaces=false
}

%==================================================
% ASSET PATH
%==================================================

\graphicspath{
    {assets/images/}
        {assets/files/}
        {assets/references/}
        {logbook/}
}

% ToC styling
\setlength{\cftbeforechapskip}{2pt}
\renewcommand{\cftchappresnum}{BAB }            % Set agar chapter punya awalan "BAB"
\renewcommand{\cftchapnumwidth}{3em}            % Tambah indentation agar tidak tabrakan
\renewcommand{\cftsecindent}{3em}               % Tambah indentation agar tidak tabrakan
\renewcommand{\cftchapfont}{\normalfont}        % Set agar chapter di daftar isi tidak bold
\renewcommand{\cftchappagefont}{\normalfont}    % Set agar numbering chapter tidak bold
\renewcommand{\cftdotsep}{0.75}                 % Set agar titik-titik tidak terlalu renggang
\renewcommand{\cftchapdotsep}{\cftdotsep}       % Set agar chapter juga punya titik-titik

\renewcommand{\cftfigindent}{0em}               % Set lof indent
\renewcommand{\cftfigpresnum}{Gambar }          % Set lof leading
\renewcommand{\cftfignumwidth}{5.5em}           % Set leading width\


% Daftar Hadir Bimbingan Table
\newenvironment{dhbtable}
{
    \begin{longtable}{| >{\centering\arraybackslash}m{1cm} | >{\justifying\arraybackslash}m{7cm} | >{\centering\arraybackslash}m{3cm} | >{\centering\arraybackslash}m{2cm} |} \hline
        \rule{0pt}{3ex} \textbf{No} &
        \rule{0pt}{3ex}\centering\textbf{Materi Bimbingan} &
        \rule{0pt}{3ex}\textbf{Tgl. Bimb.} &
        \rule{0pt}{3ex}\textbf{Paraf} \\ \hline
        \endfirsthead

        \hline
        \rule{0pt}{3ex}\textbf{No} &
        \textbf{Materi Bimbingan} &
        \textbf{Tgl. Bimb.} &
        \textbf{Paraf} 
        \endhead
}
{
    \end{longtable}
    \begin{flushright}
    {\city}, \textcolor{red}{\signaturedate}\\
    Mengetahui,\\
    Pembimbing {\jenisMBKM},\\[2cm]
    {\dosen}\\
\end{flushright}
}

% For dhbtable (dhbentry) -> insert new row
\newcommand{\dhbentry}[4]{
    #1 & 
    \fontsize{11}{14}\selectfont{#2} & 
    #3 & 
    #4
\\ \hline
}
\begin{document}
\pagenumbering{gobble}
% \renewcommand{\signaturedate}{15 April 2026}
\flushright{\small\textbf{FR-FTI-02-09/R1}}\\ % chktex 8
\begin{center}
    \large\textbf{DAFTAR HADIR \\ BIMBINGAN MAGANG INDUSTRI}\\
\end{center}
\raggedright%
Program Studi               \tab: {\studyprogram} \\
Semester/Tahun Akademik     \tab: {\semester} / {\academicperiod} \\
NIM/Nama Mahasiswa          \tab: {\authornim} / {\authorname} \\

\begin{singlespace}
    \begin{dhbtable}
        \dhbentry{1.}
        {
            Bimbingan berlangsung pada pukul 12:00--12:33 WIB % chktex 8
            secara daring menggunakan Teams Meeting. Dosen
            meminta setiap mahasiswa menjelaskan perusahaan
            tempat magang, posisi/jabatan, jobdesk yang
            dikerjakan dan framework yang digunakan.

            Dosen pembimbing memberikan arahan agar seluruh aktivitas kerja
            didokumentasikan dengan baik untuk keperluan laporan. Setiap jabatan perlu
            dijelaskan secara teoritis dan dibandingkan dengan praktik yang dilakukan
            selama magang. Mahasiswa juga diminta segera melengkapi administrasi MBKM serta
            mempersiapkan diri untuk seminar hasil yang akan dilaksanakan secara offline di
            kampus. }{Selasa,\linebreak 24 Februari 2026}{}

        \dhbentry{2.}
        {
            Bimbingan berlangsung pada pukul 19:00--20:18 WIB secara daring menggunakan
            Teams Meeting. Dosen menanyakan \textit{progress} tiap mahasiswa dan menanyakan
            terkait pengisian \textit{logbook} serta kendala yang dihadapi serta solusi
            yang dibuat mahasiswa.

            Dosen pembimbing memberikan pengarahan untuk memasukkan 3-4 foto dokumentasi
            pekerjaan ke dalam \textit{logbook}/laporan. Lalu menjelaskan proyek yang
            dikerjakan, seperti definisi, tujuan dan manfaatnya.}{Rabu,\linebreak 11 Maret
            2026}{ }

        \dhbentry{3.} { Bimbingan berlangsung pada pukul 19:00--20:20 WIB secara daring
            menggunakan Teams Meeting. Dosen memberikan evaluasi terhadap isi laporan,
            terutama pada bagian tujuan yang diminta untuk difokuskan pada tugas selama
            magang serta beberapa poin yang serupa perlu digabungkan. Dosen juga
            mengarahkan penggunaan istilah agar sesuai pedoman terbaru serta mengurangi
            penggunaan kata ganti orang pertama.% chktex 8

            Dosen pembimbing menegaskan bahwa dokumen informal seperti tangkapan layar
            percakapan dipindahkan ke lampiran, sementara sumber formal tetap dapat
            digunakan dalam laporan. Selain itu, mahasiswa diminta menambahkan dokumentasi
            kegiatan pada logbook sebanyak 3-4 foto per bulan, menjelaskan alasan pemilihan
            bidang pada latar belakang, serta memastikan format penulisan daftar pustaka
            yang digunakan sudah sesuai.}{Rabu,\linebreak 1 April 2026}{}

        \dhbentry{4.} {Finalisasi
            Laporan Kemajuan Magang dan penandatanganan berkas-berkas}{Kamis,\linebreak 16 April
            2026}{}%chktex 8

        \dhbentry{5.} {
            Berlangsung pada pukul 12:01--12:24 WIB secara daring menggunakan Teams
            Meeting. Dosen memberikan arahan terkait bagian-bagian yang harus diisi
            dalam Laporan Akhir Magang Industri, terutama pada Bab 4 dan Bab 5.
            Selain itu, dosen juga menjelaskan poin-poin penting yang perlu
            diperhatikan agar isi laporan dapat disusun secara sistematis dan sesuai
            dengan ketentuan yang berlaku.
        }{Selasa,\linebreak 19 Mei 2026}{}

        \dhbentry{6.} {
            Berlangsung pada pukul 19:00--19:30 WIB secara daring menggunakan Teams
            Meeting. Dosen memberikan koreksi pada laporan masing-masing mahasiswa,
            terutama terkait format penulisan, konsistensi penyajian, dan kejelasan
            isi laporan. Masukan yang diberikan digunakan sebagai acuan untuk
            melakukan revisi sehingga laporan menjadi lebih terstruktur dan mudah
            dipahami.
        }{Kamis,\linebreak 11 Juni 2026}{}

        \dhbentry{7.} {
            Berlangsung pada pukul 12:00--13:00 WIB secara luring di Ruang Rapat,
            Lantai 11 Gedung R, Kampus 1 Universitas Tarumanagara. Dosen memberikan
            koreksi terhadap \textit{powerpoint} yang akan digunakan pada Sidang
            Magang, khususnya terkait susunan materi dan kejelasan penyampaian.
            Selain itu, dilakukan pula evaluasi terhadap kelengkapan berkas-berkas
            yang harus dikumpulkan menjelang pelaksanaan Sidang Magang guna memastikan
            seluruh persyaratan telah terpenuhi.
        }{Senin,\linebreak 15 Juni 2026}{}

    \end{dhbtable}
\end{singlespace}
\end{document}